% Architecture & Stack -- BNP x Mistral assistant
% Compile with: lualatex report.tex   (or xelatex report.tex)
% Reason: source contains UTF-8 glyphs (French quotes, accents, math, emoji).
\documentclass[11pt]{article}

\usepackage[margin=1in]{geometry}
\usepackage{fontspec}
\usepackage{polyglossia}
\setdefaultlanguage{french}
\setotherlanguage{english}
\usepackage{listings}
\usepackage{xcolor}
\usepackage{booktabs}
\usepackage{longtable}
\usepackage{array}
\usepackage{enumitem}
\usepackage{hyperref}
\usepackage{amssymb}
\usepackage{eurosym}

\hypersetup{
  colorlinks=true,
  linkcolor=blue!50!black,
  urlcolor=blue!50!black,
}

\definecolor{codebg}{RGB}{246,246,246}
\definecolor{codeframe}{RGB}{210,210,210}
\definecolor{kw}{RGB}{20,80,160}
\definecolor{str}{RGB}{40,120,40}
\definecolor{cmt}{RGB}{120,120,120}

\lstdefinestyle{py}{
  language=Python,
  basicstyle=\ttfamily\small,
  keywordstyle=\color{kw}\bfseries,
  stringstyle=\color{str},
  commentstyle=\color{cmt}\itshape,
  showstringspaces=false,
  breaklines=true,
  breakatwhitespace=true,
  columns=fullflexible,
  keepspaces=true,
  frame=single,
  rulecolor=\color{codeframe},
  backgroundcolor=\color{codebg},
  numbers=none,
  xleftmargin=4pt,
  xrightmargin=4pt,
  literate=
    {é}{{\'e}}1 {è}{{\`e}}1 {ê}{{\^e}}1 {à}{{\`a}}1 {â}{{\^a}}1
    {ô}{{\^o}}1 {ç}{{\c{c}}}1 {î}{{\^i}}1 {ï}{{\"i}}1 {û}{{\^u}}1
    {«}{{\guillemotleft{}}}1 {»}{{\guillemotright{}}}1 {€}{{\euro{}}}1,
}

\lstdefinestyle{plain}{
  basicstyle=\ttfamily\scriptsize,
  breaklines=false,
  columns=fullflexible,
  keepspaces=true,
  frame=single,
  rulecolor=\color{codeframe},
  backgroundcolor=\color{codebg},
  xleftmargin=4pt,
  xrightmargin=4pt,
}

\title{Architecture \& Stack --- BNP $\times$ Mistral assistant}
\author{Matthieu Minguet}
\date{}

\begin{document}
\maketitle

\section{Use case}

24/7 conversational assistant for a retail-banking client that:

\begin{enumerate}[leftmargin=*]
  \item \textbf{Answers} broad customer-service questions --- balance, transactions, contract terms, agency hours, advisor contact.
  \item \textbf{Acts} on low-risk reversible operations --- internal transfers between the client's own accounts, temporary card lock.
  \item \textbf{Escalates} anything risky or irreversible --- outbound wires, profile changes --- to a human advisor. Enforced at the \textbf{system layer}, not just the prompt.
\end{enumerate}

The prototype models a single client (me) with realistic BNP-style products: Compte de chèques, Livret A, portefeuille investissement, Visa Premier, prêt études, assurance auto, et le contrat Esprit Libre.

\pagebreak
\section{Stack at a glance}

\begin{lstlisting}[style=plain]
+------------------------------------------------------------------+
| Browser (vanilla HTML/CSS/JS, BNP theme)                         |  static/index.html
|   * SSE consumer . streaming markdown bubbles                    |  static/app.js
|   * tool-call chips . blocked indicator                          |  static/style.css
|   * LIVE accounts/cards/products sidebar (polls /state)          |
+----------------------------+-------------------------------------+
                             | HTTP . text/event-stream
+----------------------------v-------------------------------------+
| FastAPI                                                          |  app/main.py
|   GET  /         -> static/index.html                            |
|   GET  /state    -> live JSON snapshot of mock bank              |
|   POST /chat     -> SSE stream from agent loop                   |
+----------------------------+-------------------------------------+
                             |
+----------------------------v-------------------------------------+
| Agent loop                                                       |  app/agent.py
|   * Ollama chat client (streaming + native tool-calling)         |
|   * normalize tool calls -> guardrails -> tool -> re-prompt      |
|   * Hard cap: 5 tool iterations / turn                           |
+--------+-------------------+-------------------------------------+
         |                   |
   +-----v-----+       +-----v--------------------------+
   | Ollama    |       | Tool dispatcher                |  app/tools.py
   | (local)   |       |   * arg coercion + filtering   |
   | ministral |       |   * registry of 12 tools       |
   |   -3:8b   |       |   * routes through guardrails  |  app/guardrails.py
   +-----------+       +-------+------------+-----------+
                               |            |
                       +-------v------+  +--v-------------+
                       | SQLite       |  | Chroma RAG     |  app/rag.py
                       | mock bank    |  |  hybrid scoring|
                       | (5 tables)   |  | over .md docs  |  data/contracts/*.md
                       +--------------+  +----------------+
                       app/db.py
                       data/bank.sqlite
\end{lstlisting}

\begin{longtable}{@{}p{2.6cm}p{4.4cm}p{7.5cm}@{}}
\toprule
\textbf{Layer} & \textbf{Choice} & \textbf{One-line justification} \\
\midrule
\endhead
Inference & Ollama, local & Data sovereignty, GDPR posture, no token cost, no model leakage \\
Model & \texttt{ministral-3:8b} (default) / \texttt{:14b} & Mistral family, native function-calling. 8B fits $\sim$6\,GB VRAM, 14B $\sim$10\,GB \\
Orchestration & Plain Python loop & Easier to demo and audit than LangChain/LangGraph for a prototype \\
Tool wiring & In-process Python & One process, no IPC; can be swapped for tooling later \\
Bank data & SQLite & File-based, zero setup, easy to reset and seed \\
RAG store & Chroma (local persistent) & No server, no network call, embeds \& queries from disk \\
Embeddings & \texttt{nomic-embed-text} (Ollama) & Same local-only story as the LLM; multilingual \\
API & FastAPI + SSE (\texttt{sse-starlette}) & Native streaming, type-safe, minimal boilerplate \\
Frontend & Vanilla HTML/JS + \texttt{marked} (CDN) & No build step, transparent for a panel walkthrough \\
\bottomrule
\end{longtable}

\section{Why local Ollama + Ministral}

Simulates running the model on customer hardware, which is needed for three reasons:

\begin{enumerate}[leftmargin=*]
  \item \textbf{Regulatory / GDPR.} A retail-bank assistant handles position data, transaction labels, IBANs. Routing that to a hosted endpoint either crosses borders or requires a complex DPA. Local inference keeps personal data on the bank's infrastructure. Furthermore, personally it is a way for me to experiment with the product easily and cheaply.
  \item \textbf{No training feedback loop.} Anything sent to a hosted provider is governed by their retention policy. A local model has zero risk of contributing customer data back into training. It also answers to GRPD requirements.
  \item \textbf{Latency control.} TTFT depends on the network only when the model is hosted. Locally we control the entire path and can expect a speed-up on profesional hardware; on a consumer GPU \texttt{ministral-3:8b} lands well under the 2s ceiling and exceeds reading speed.
\end{enumerate}

Trade-off: an 8B local model is weaker at deep reasoning than a frontier hosted model, and sometimes misses the correct syntax for the tool use. We mitigate by \textbf{forcing tool use for all factual answers}, meaning the model orchestrates, the tools provide ground truth, and a robust dispatcher (next section) absorbs the model's imperfect arg shapes.

Plus corresponds to what could be done by the bank on-premises given its infrastructure and regulatory constraints.

\section{The agent loop}

\texttt{app/agent.py:run\_turn} is a generator. It streams from Ollama, intercepts emitted tool calls, runs them through guardrails + tools, feeds results back, and loops.

\begin{lstlisting}[style=py]
def run_turn(history):
    messages = [{"role": "system", "content": SYSTEM_PROMPT}, *history]

    for _ in range(MAX_TOOL_ITERATIONS):                   # cap = 5
        stream = ollama.chat(MODEL, messages=messages,
                             tools=tools.SCHEMAS, stream=True)
        accumulated_content, accumulated_tool_calls = "", []
        for chunk in stream:
            msg = chunk["message"] or chunk.message
            if msg.content:
                accumulated_content += msg.content
                yield {"type": "token", "text": msg.content}
            if msg.tool_calls:
                accumulated_tool_calls.extend(msg.tool_calls)

        messages.append({"role": "assistant", ...})
        if not accumulated_tool_calls:
            yield {"type": "done"}; return

        for tc in accumulated_tool_calls:
            name, args = _normalize_tool_call(tc)
            yield {"type": "tool_start", "name": name, "args": args}
            allowed, result = tools.run_tool(name, args)
            yield {"type": "tool_result", "name": name,
                   "allowed": allowed, "result": result}
            messages.append({"role": "tool", "name": name,
                             "content": json.dumps(result)})
\end{lstlisting}

Why streaming: TTFT is what the user perceives. Streaming makes the first token visible the moment Ollama produces it.

Why \texttt{\_normalize\_tool\_call}: the Ollama Python client emits tool-call shapes that vary across versions and model behaviours --- sometimes \texttt{arguments} is a dict, sometimes a JSON-encoded string, sometimes empty/malformed. Normalizing in one place keeps the rest of the loop trivial:

\begin{lstlisting}[style=py]
def _normalize_tool_call(tc):
    fn = tc["function"] if isinstance(tc, dict) else tc.function
    name = fn["name"] if isinstance(fn, dict) else fn.name
    raw_args = fn["arguments"] if isinstance(fn, dict) else fn.arguments
    if isinstance(raw_args, str):
        try: args = json.loads(raw_args) if raw_args else {}
        except json.JSONDecodeError: args = {}
    elif isinstance(raw_args, dict):
        args = raw_args
    else:
        args = dict(raw_args)
    return name, args
\end{lstlisting}

Hard cap: a misaligned model can bounce between tools forever. 5--20 iterations covers any realistic chain (RAG $\rightarrow$ balance read $\rightarrow$ action) and fails loudly if exceeded, while staying quick enough.

\section{Tool dispatcher --- robustness as a feature}

A small open-weight model gets argument shapes wrong often, which is something I found from testing: extra keys, numbers as strings, French decimals (\texttt{"100,50"}), boolean strings (\texttt{"oui"}), unknown account names. If the dispatcher is brittle, the agent loop dies and the user sees a failure message.

\texttt{app/tools.py:run\_tool} is the only path to tool execution. It does three things, in order:

\begin{lstlisting}[style=py]
def run_tool(name, args):
    fn = REGISTRY.get(name)
    if fn is None:
        return False, {"error": f"unknown tool '{name}'"}
    if not isinstance(args, dict):
        args = {}

    clean_args = _normalize_args(fn, args)        # 1. filter + coerce
    decision = guardrails.check(name, clean_args) # 2. data-coded rules
    if not decision.allowed:
        return False, {"blocked": True, "reason": decision.reason}

    try:
        return True, fn(**clean_args)             # 3. invoke
    except TypeError as e:
        return False, {"error": f"bad arguments: {e}"}
    except ValueError as e:
        return False, {"error": str(e)}
\end{lstlisting}

\texttt{\_normalize\_args} reads each tool's signature, drops unknown kwargs silently, and coerces by annotation:

\begin{lstlisting}[style=py]
def _normalize_args(fn, raw):
    sig = inspect.signature(fn)
    out = {}
    for pname, param in sig.parameters.items():
        if pname not in raw: continue
        v = raw[pname]
        target = _ANN_MAP.get(param.annotation)   # handles `from __future__`
        if target is not None:                    # -> annotations are strings
            try: v = _coerce(v, target)
            except ValueError: pass
        out[pname] = v
    return out
\end{lstlisting}

\texttt{\_coerce} accepts the shapes a real Mistral 8B has been observed to emit:

\begin{itemize}[leftmargin=*]
  \item \texttt{float}: \texttt{100}, \texttt{100.0}, \texttt{"100"}, \texttt{"100.50"}, \texttt{"100,50"} (French decimal)
  \item \texttt{int}: \texttt{"3"}, \texttt{4.0}, \texttt{True}
  \item \texttt{bool}: \texttt{True}, \texttt{"true"}/\texttt{"True"}, \texttt{"yes"}/\texttt{"oui"}, \texttt{"1"}, \texttt{1} (and falsy mirrors)
\end{itemize}

The 12 tools exposed to the model (see \texttt{tools.SCHEMAS}):

\begin{longtable}{@{}lll@{}}
\toprule
\textbf{Read} & \textbf{Write} & \textbf{RAG / Reference} \\
\midrule
\endhead
\texttt{get\_client\_profile} & \texttt{transfer\_internal} & \texttt{search\_contracts} \\
\texttt{list\_accounts} & \texttt{lock\_card} & \texttt{get\_counselor\_contact} \\
\texttt{get\_account\_balance} & \texttt{transfer\_external} \emph{(always blocked)} & \texttt{get\_agency\_hours} \\
\texttt{list\_transactions} & & \\
\texttt{list\_cards} & & \\
\texttt{list\_products} & & \\
\bottomrule
\end{longtable}

\texttt{get\_client\_profile} is the entry point: the system prompt instructs the model to call it whenever it needs the client's name, owned accounts, or product list. The agent never has identity hardcoded.

\section{Mock banking layer}

\texttt{app/db.py} is a thin SQLite layer over JSON fixtures (\texttt{data/fixtures/*.json}). The agent doesn't know it's mocked --- it sees the tool API. To swap in a real backend, replace the bodies of \texttt{db.list\_accounts}, \texttt{db.get\_transactions}, etc. The tool schemas in \texttt{app/tools.py:SCHEMAS} stay identical, the guardrails stay identical, the system prompt stays identical.

In a production BNP context, the seams would map to:

\begin{itemize}[leftmargin=*]
  \item \textbf{Read tools} $\rightarrow$ existing core-banking system / Cobol mainframe APIs, typically already exposed via an internal REST gateway.
  \item \textbf{Action tools} $\rightarrow$ the real movement engine and card-management system, with idempotency keys and audit logging.
  \item \textbf{Reference data} $\rightarrow$ the CRM.
\end{itemize}

Since the tool layer is the contract, swapping in MCP later is a refactor --- each of the 12 tools becomes one MCP \texttt{tool} definition, and human-escalation routing becomes a second MCP server querying the CRM by topic / availability.

\section{RAG design --- hybrid retrieval}

\textbf{Why per-product contracts.} The corpus mirrors the products the client owns (six markdown files). Index small, retrieval precise, citations meaningful. In production, gate retrieval by what the authenticated client is actually subscribed to --- never let a model surface a clause for a product the client doesn't have.

\textbf{Pipeline:}

\begin{enumerate}[leftmargin=*]
  \item \textbf{Ingestion} (\texttt{rag.py:build\_index}) --- paragraph-aware chunking, $\sim$500 chars with 80-char overlap. Source filename embedded in the chunk text (\texttt{[esprit\_libre] \# Découvert autorisé \ldots}) so the embedding picks up document identity.
  \item \textbf{Embeddings} --- \texttt{nomic-embed-text} via Ollama (multilingual, $\sim$270\,MB).
  \item \textbf{Store} --- Chroma persistent client at \texttt{data/chroma/}, cosine distance.
\end{enumerate}

\begin{lstlisting}[style=py]
def retrieve(query, k=4):
    if not query.strip(): return []
    qtoks = _tokens(query)

    # Vector candidates (top 3k)
    res = coll.query(query_embeddings=[_embed([query])[0]],
                     n_results=max(k * 3, 12))
    seen = {f"{m['source']}#{m['chunk']}": {...} for ...}

    # Lexical pass: pull ALL chunks whose filename matches a query token.
    # Catches "Esprit Libre", "Livret A" when vector recall buries them.
    for doc, meta in coll.get(...):
        if any(t in meta["source"].lower() for t in qtoks):
            seen.setdefault(...)

    # Combine: vec_score + 0.05 * token_overlap + 0.6 * src_match
    candidates = [...]
    candidates.sort(key=lambda c: c["score"], reverse=True)
    return candidates[:k]
\end{lstlisting}

The agent prompt requires it to cite sources from \texttt{search\_contracts} (e.g.\ «~selon \texttt{esprit\_libre.md}~»).

In production, replace Chroma with the bank's vector index of choice (pgvector, OpenSearch), gate by client identity, and add a learned re-ranker if the latency budget permits. I currently don't have all of that information available nor the need for it in a prototype.

\section{Guardrails --- three layers, code-coded rules}

I tried to implement multiple fail safes for sensitive operations, especially \texttt{transfer\_external} which simulates an outbound wire and represents an irreversible action with fraud risk.

\begin{enumerate}[leftmargin=*]
  \item \textbf{Prompt layer} --- the system prompt tells the model what's blocked and instructs it not to bypass.
  \item \textbf{Schema layer} --- \texttt{transfer\_external} exists in the tool schema (so the model can attempt the call and learn it fails) but its dispatcher always returns a structured \texttt{blocked} payload.
  \item \textbf{Code layer} --- \texttt{app/guardrails.py:check} is the only path through \texttt{run\_tool}.
\end{enumerate}

\begin{lstlisting}[style=py]
def check(tool_name, args) -> Decision:
    if tool_name == "transfer_external":
        return Decision(False, "Les virements externes (...) sont bloqués...")

    if tool_name == "transfer_internal":
        owned = _owned_account_ids()
        from_id, to_id = args.get("from_account"), args.get("to_account")
        try: amount = float(args.get("amount", 0))
        except (TypeError, ValueError):
            return Decision(False, "Le montant fourni n'est pas un nombre valide.")
        if from_id not in owned or to_id not in owned:
            return Decision(False, "Compte source ou destination inconnu.")
        if from_id == to_id:
            return Decision(False, "Compte source et destinataire identiques.")
        if amount <= 0:
            return Decision(False, "Le montant doit être strictement positif.")
        if amount > 10_000 and not args.get("confirmed"):
            return Decision(False, "Montant > 10 000 €. Confirmation requise.")
        src = db.get_account(from_id)
        if src and src["balance"] - amount < 0:
            return Decision(False, f"Solde insuffisant ({src['balance']:.2f} €). "
                                   "Découvert non autorisé.")
        return Decision(True)

    if tool_name in {"get_account_balance", "list_transactions"}:
        if (acc := args.get("account_id")) and acc not in _owned_account_ids():
            return Decision(False, "Ce compte n'appartient pas au client.")
        return Decision(True)

    return Decision(True)
\end{lstlisting}

\begin{longtable}{@{}p{8.5cm}p{6cm}@{}}
\toprule
\textbf{Rule} & \textbf{Effect} \\
\midrule
\endhead
\texttt{transfer\_external} & Always blocked. Routes user to advisor / strong auth. \\
\texttt{transfer\_internal} $\land$ unknown account & Blocked. \\
\texttt{transfer\_internal} $\land$ same source/destination & Blocked. \\
\texttt{transfer\_internal} $\land$ amount $\leq 0$ & Blocked. \\
\texttt{transfer\_internal} $\land$ amount $>$ 10\,000\,\euro{} $\land$ $\lnot$ confirmed & Blocked. \\
\texttt{transfer\_internal} $\land$ source balance $-$ amount $<$ 0 & Blocked (no overdraft). \\
Account id not owned by the client & Blocked. \\
All read tools, \texttt{lock\_card}, \texttt{search\_contracts}, \texttt{get\_*} & Allowed. \\
\bottomrule
\end{longtable}

When a guardrail fires, the UI surfaces a red \texttt{[bloqué: <tool>]} chip on the assistant message.

The 10\,000\,\euro{} confirm gate is intentionally below the overdraft check: a 50\,000\,\euro{} transfer with \texttt{confirmed=true} is still blocked on overdraft. Two independent gates, both must pass.\\

Regarding information, the model is loaded for the client, to avoid it leaking contracts from other products or clients. However, the model is not fine-tuned on the client's data, which it has to learn to use the tools and RAG to access that data.

\section{Latency budget}

Targets for the product:

\begin{itemize}[leftmargin=*]
  \item \textbf{TTFT $<$ 2\,s} --- perceived first-token latency.
  \item \textbf{Tokens/sec $>$ reading speed} --- roughly $>$ 5 tok/s.
\end{itemize}

What contributes to TTFT, in order:

\begin{enumerate}[leftmargin=*]
  \item Model load (cold) --- Ollama keeps the model resident after first request; first call after boot pays a one-time cost.
  \item Prompt prefix --- system prompt + 12 tool schemas + history. Caches well in Ollama once warm.
  \item First inference token --- depends on GPU (with my 5090 this is very quick).
\end{enumerate}

Measured on this machine: a one-tool turn (\textit{«~Solde Livret A ?~»}) returns the full streamed response over HTTP in \textbf{$\sim$290\,ms} warm.

\section{What changes in production}

\begin{itemize}[leftmargin=*]
  \item \textbf{Authentication \& authorization.} Today: single hardcoded client via \texttt{get\_client\_profile}. Real: SSO / mobile auth, per-tool scopes, session timeouts, step-up auth for any write action. Could be linked to the current identity check when calling through phone, or identification through the app.
  \item \textbf{Real database.} Today: SQLite with JSON fixtures. Real: the bank's core banking system, queried through an internal API layer.
  \item \textbf{Legacy system integration.} The bank has existing infrastructure for transactions, card management, CRM, etc, often in old languages like Cobol and need to integrate with them.
  \item \textbf{Multi-tenant data isolation.} Today: one SQLite. Real: per-client row-level security, every tool scoped by authenticated client id.
  \item \textbf{Audit log.} Every tool call (allowed \emph{and} blocked) persisted with user id, args, timestamp, model id, prompt hash. Required for regulatory traceability the agent loop already emits these as events, so just need a sink.
  \item \textbf{Risk-model integration.} Hook the bank's existing fraud / risk engine into \texttt{guardrails.check()} and see even if new models could be developped.
  \item \textbf{Observability.} Latency histograms (TTFT, tok/s, tool-call duration), cost-equivalent token counts, refusal rate, guardrail-fire rate.
  \item \textbf{Eval harness.} Golden set of $\sim$100 user queries with expected tool calls and answer rubrics; regress on every model upgrade. Plus needs to pass the ``grandma test'': would my grandma be able to use it without it doing something wrong?
  \item \textbf{Human escalation routing.} The brainstorm's routing idea: a second tool/server that, given the user's topic and frustration signal, picks the right advisor or queue.
  \item \textbf{Compliance review.} Any RAG citation that contradicts a current legal document is a liability --- versioning of the contract corpus, signed snapshots.
  \item \textbf{UX.} Confirmation modals for any write action, transcript export, language toggle, accessibility, dark mode.
\end{itemize}

\section{Quick demo script}

Created a \texttt{./run.sh} that launches everything (Ollama daemon if needed, missing model pulls, venv/conda setup, seed, uvicorn). Then in the browser some of the things able to be tested are:

\begin{enumerate}[leftmargin=*]
  \item \textit{«~Liste mes comptes.~»} --- tools: \texttt{get\_client\_profile} $\rightarrow$ \texttt{list\_accounts}.
  \item \textit{«~Quel est le solde de mon Livret A ?~»} --- single tool call, fast. Sidebar matches.
  \item \textit{«~Mes 5 dernières transactions sur le compte de chèques.~»} --- formatted markdown list.
  \item \textit{«~Que dit mon contrat Esprit Libre sur le découvert autorisé ?~»} --- RAG hit, citation \texttt{esprit\_libre.md}.
  \item \textit{«~Vire 1000~\euro{} de mon compte de chèques vers mon Livret A.~»} --- write executes, \textbf{sidebar flashes green and updates}.
  \item \textit{«~Vire 50~000~\euro{} de mon compte de chèques vers mon Livret A.~»} --- guardrail blocks (overdraft), red \texttt{[bloqué]} chip, sidebar unchanged.
  \item \textit{«~Envoie 5000~\euro{} à l'IBAN FR76\ldots~»} --- guardrail blocks (external), red chip, model pivots to advisor.
  \item \textit{«~Bloque ma carte Visa Premier.~»} --- reversible action, card pill flips active $\rightarrow$ Bloquée.
  \item \textit{«~Quels sont les horaires de mon agence le samedi ?~»} --- static lookup.
\end{enumerate}

\end{document}
